
\section{模型的评价}

\subsection{模型的优点}
\begin{itemize}[itemindent=2em]
\item \textbf{优点1：}几何模型精确且考虑物理细节。本文建立了三维线段间最短距离的解析算法，严格处理了周期性边界条件，通过27个周期镜像的枚举确保了距离计算在环面拓扑下的精确性，避免了几何近似带来的系统误差。
\item \textbf{优点2：}高效的连通性判定算法。引入空间哈希网格将颗粒间距离计算的复杂度从$O(N^2)$降至实际接近$O(N)$，并采用并查集数据结构支持动态连通性判定和提前终止，使得700颗粒规模下的蒙特卡洛模拟可在秒级完成单次试验。
\item \textbf{优点3：}系统的统计推断方法。从单颗粒确定性分析到多颗粒蒙特卡洛模拟，再到二分搜索优化，构成了完整的统计建模链条。每次概率估计均报告标准误差和置信区间，确保了结论的统计严谨性。
\item \textbf{优点4：}揭示了高长径比填料的超低渗透阈值现象。通过系统的数值实验，定量确定了长径比83:1的圆柱形填料在体积分数仅为0.0551\%时即可实现90\%以上导通概率，为导电复合材料设计提供了明确的定量指导。
\end{itemize}

\subsection{模型的缺点}
\begin{itemize}[itemindent=2em]
\item \textbf{缺点1：}忽略了颗粒间可能的空间排斥效应。本文假设颗粒位置完全随机且独立，未考虑实际制造过程中颗粒间的体积排斥（硬核势）和可能的团聚效应。在较高填充率下，体积排斥可能导致颗粒分布偏离完全随机假设。
\item \textbf{缺点2：}球体-圆柱体距离计算存在近似。点（球心）到线段（圆柱体轴线）的距离计算使用了周期性枚举法，在球体靠近微构体面、角或边的极端情况下可能存在微小偏差，影响连通判定的精度。
\end{itemize}
